\documentclass[12pt]{article}
\usepackage{float}
\usepackage{graphicx} % Required for inserting images
\usepackage{fontspec}
\usepackage{apacite}
\setmainfont{Times New Roman}
\usepackage[letterpaper,top=2cm,bottom=2cm,left=2cm,right=2cm,marginparwidth=1.5cm]{geometry}

\title{Weekly Assignment 2: Reflection on Cybersecurity Culture}
\author{Truong Khang Tobias Nguyen}
\date{}

\begin{document}

\maketitle

\section{Cybersecurity Culture in Norway}

Based on my analysis, Norway's cybersecurity culture is based on trust, collectivism and resilience according to national policy and security guidance.
\newline
While the national security strategy states that security is a societal responsibility, the Norwegian National Security Authority states that every organization and individual has to contribute their part to ensure national resilience \cite{natsecstrat, nsm}. 
\newline
Similarly, the Norwegian Center for Information Security (NorSIS) describes Norway as digitally interconnected, meaning that attacks on an individual or an organization can have effects beyond the original victim \cite{norsis}. 
\newline
One assumption that can be made from these reports is that security is a shared responsibility where every individual in a system contributes to the security of their systems and organizations.

\section{Reflections}
Overall, my assumptions on cybersecurity aligns with the Norwegian approach that security has socio-technical nature rather than a purely technical one.
\newline
Although I largely support the assumption that security should be a shared responsiblity, I am sceptical on putting the responsibility across individuals. Secure behaviour is not only about how individuals act but also about organizational incentives, system architecture and leadership. In a sense, I see individuals as one layer of cybersecurity, but not the entire cybersecurity architecture itself.

\section{Country-wide Culture and their Effects on Individual Behaviour}
I do believe that a country-wide culture affects individual behaviour, but it does not determine it. A country-wide culture provides individuals with a certain baseline about trust, responsibility and behaviour which can affect an organisation in the way it operates. 
\newline
Nevertheless, I also think that organizations can also either reinforce or contradict country-wide culture. I believe that that security culture is shaped by the individual organisation and it's leadership. 
\newline
Thus, an employee may come from a society that values cooperation and trust, but whether they report an incident or bypass a security measure dependd on what their organisation rewards or discourages.

\bibliographystyle{apacite}
\bibliography{references}

\end{document}